\documentclass[12pt]{article}
\usepackage[utf8]{inputenc}
\usepackage[T1]{fontenc}
\usepackage[margin=1in]{geometry}
\usepackage{setspace}
\usepackage{xcolor}
\usepackage[colorlinks=true, urlcolor=blue!60!black, linkcolor=blue!60!black]{hyperref}
\setlength{\parindent}{0em}
\setlength{\parskip}{0.45em}
\renewcommand{\baselinestretch}{1.0}
\begin{document}
\begin{center}
{\Large \textit{Students for Fair Admissions, Inc. v. President and Fellows of Harvard College}, 600 U.S. 181 (2023)\par}
\vspace{0.3em}
{\normalsize Case Brief --- One-Page Edition\par}
\end{center}

\textbf{Memory Jogger.} The end of race-conscious college admissions (\href{https://www.law.cornell.edu/supremecourt/text/20-1199}{600 U.S.\ 181}).

\textbf{Facts.} Students for Fair Admissions sued Harvard and UNC over holistic admissions that considered race, including Harvard's tracking of its class's racial makeup and UNC's racial ``plus.''

\textbf{Procedural History.} After bench trials, both district courts upheld the programs; the First Circuit affirmed in \textit{Harvard}.

\textbf{Votes.} Reversed, 6--2 and 6--3 (Roberts for the Court; Thomas, Gorsuch, and Kavanaugh concurring; Sotomayor and Jackson dissenting).

\textbf{Holding.} Both programs fail strict scrutiny under the Equal Protection Clause and Title VI: unmeasurable goals, race used as a negative and a stereotype, and no logical endpoint. Applicants may still discuss how race affected their lives, tied to their character or ability. Military academies were reserved.

\textbf{Reasoning.} The Clause is universal; diversity goals were ``commendable'' but unreviewable, racial categories imprecise, preferences zero-sum, and racial shares stable year after year. Thomas defended colorblind original meaning; Gorsuch relied on Title VI's text. Sotomayor accused the Court of ``cement[ing] a superficial rule of colorblindness''; Jackson documented ``gulf-sized race-based gaps.''

\textbf{Examples.} \emph{Same-Side:} a medical school's fixed racial bands. \emph{Opposite-Side:} credit for an essay about overcoming discrimination. \emph{Fence-Sitter:} military academies.

\textbf{Critique.} Originalists praise anti-classification but question the military carve-out; critics invoke anti-subordination history.

\textbf{Key Quotation.} ``Eliminating racial discrimination means eliminating all of it. And the Equal Protection Clause, we have accordingly held, applies `without regard to any differences of race, of color, or of nationality'---it is `universal in [its] application.' '' --- Chief Justice Roberts, 600 U.S.\ at 206.

\end{document}
